\section{The metric on the singular strata}\label{sec:strata}

We finish by distinguishing the quadratic approximation
from an exact local cone metric. The argument uses only
the intrinsic recognition of the two singular strata and
their harmonic metrics.

\begin{lemma}\label{lem:intrinsic-strata}
The subset $\Sigma$ consists precisely of the points of
$M$ whose metric tangent cones are not Euclidean. Within
$\Sigma$, the subset $\Delta$ is its topological
nonmanifold locus. The corresponding descriptions hold
in the cones \eqref{eq:offdiag-cone} and
\eqref{eq:diagonal-cone}.
\end{lemma}

\begin{proof}
At a regular point the tangent cone is Euclidean.
At an off-diagonal point the normal cone is
$\mathcal N_q=T^*\CC^q\hkr_0\U(1)$. Its link is
\begin{equation}\label{eq:link}
 L_q=V_2(\CC^q)/S^1,
\end{equation}
where $V_2(\CC^q)$ is the complex Stiefel manifold
of orthonormal two-frames and the circle multiplies
both vectors by the same phase. To see this, conjugate
the cotangent coordinate in the equations
$y(x)=0$, $|x|=|y|$ and then fix the common norm.

The fibration
$S^{2q-3}\to V_2(\CC^q)\to S^{2q-1}$ gives
$\pi_1(V_2)=\pi_2(V_2)=0$ for $q\ge4$.
The circle bundle over $L_q$ consequently gives
$\pi_1(L_q)=0$ and $\pi_2(L_q)=\ZZ$. Hurewicz
then gives $H_2(L_q,\ZZ)=\ZZ$. Thus the vertex
of $\mathcal N_q$ has nonzero local homology in
degree three, strictly below its real dimension
$4q-4$. Multiplication by the Euclidean factor
shifts this nonzero local homology but does not
remove it. The off-diagonal cone is therefore not
a topological manifold at its vertex, and in
particular is not Euclidean.

The diagonal cone contains off-diagonal points
arbitrarily close to its vertex: choose two
distinct one-dimensional loop characters tending
to zero and take their direct sum. At each such
point the normal slice has the off-diagonal form
just considered. Hence the diagonal cone also
cannot be Euclidean. Uniqueness of tangent cones
from Section~\ref{sec:distances} proves the first
assertion on $M$. The finite-dimensional slice
theorem gives the same assertion at all points
of each cone.

By \eqref{eq:strata}, $\Sigma=\Sym^2N$.
Off the diagonal it is a manifold. Near a
diagonal point its germ is
\begin{equation}\label{eq:symmetric-square}
 \CC^n\times(\CC^n/\{\pm1\}).
\end{equation}
The normal link is $\mathbb{RP}^{2n-1}$, whose
first integral homology is $\ZZ/2$. This gives
nontop local homology at the vertex, so the
nonmanifold locus of $\Sigma$ is precisely
$\Delta$. In the diagonal cone, the nonsimple
semisimple representations are unordered pairs
of one-dimensional loop characters. Their
space is $\Sym^2\CC^n$, with the same description
\eqref{eq:symmetric-square}; its nonmanifold
locus is the diagonal flat factor $\CC^n$.
In the off-diagonal cone the singular set is
the flat factor $\CC^{2n}\times\{0\}$ and is
everywhere a manifold. This proves all the
claimed intrinsic descriptions.
\end{proof}

\begin{lemma}\label{lem:stratum-length}
Let $S_0$ be either $\Delta$ or
$\Sigma\setminus\Delta$. The intrinsic length
metric obtained by restricting curves to $S_0$
and measuring their lengths with the completed
ambient distance is the Riemannian length metric
of its harmonic stratum metric. The same assertion
holds for the corresponding flat strata in the
tangent cones.
\end{lemma}

\begin{proof}
Fix $s\in S_0$. In the chart centred at $s$, the
stratum is the fixed factor $U\times\{0\}$.
The coarse tensor bounds and the exit estimate
make the restricted ambient distance locally
bi-Lipschitz equivalent to a smooth metric on
$S_0$. Thus a rectifiable curve has a locally
absolutely continuous parametrization in its
smooth coordinates.

For a differentiable curve $\gamma$ with
$\gamma(0)=s$, write its fixed-factor coordinate
as $u(t)=t\dot u(0)+o(t)$. The radius formula
\eqref{eq:radius-main} in this fixed centred
chart gives
\begin{equation}\label{eq:metric-speed}
 \lim_{t\to0}\frac{D(\gamma(t),\gamma(0))}{|t|}
       =\|\dot u(0)\|_U
       =\|\dot\gamma(0)\|_{g_{S_0}}.
\end{equation}
The last equality follows from the harmonic
normalization at the centre. The argument also
applies at almost every differentiability point
of an absolutely continuous curve. It is
pointwise: the comparison constant belongs to
the fixed centre at which the limit is taken.

The length of a Lipschitz curve in a metric
space is the integral of its metric speed
\cite[Theorem~2.7.6]{bbi01}. Applying this
locally, and subdividing a rectifiable curve,
proves equality of its ambient-metric length
and harmonic Riemannian length. Infima over
curves in the stratum prove the assertion.
For the finite cones the same proof uses the
finite metric slice and its radius comparison.
\end{proof}

\begin{proposition}\label{prop:not-isometric}
For every singular point $p\in M$, no pointed
neighbourhood of $(M,D,p)$ is isometric to a
neighbourhood of the vertex in its metric
tangent cone.
\end{proposition}

\begin{proof}
An isometry preserves metric tangent cones,
and hence preserves $\Sigma$ by
Lemma~\ref{lem:intrinsic-strata}. If $p$ is
off the diagonal, a sufficiently small
neighbourhood meets only the smooth stratum
$\Sigma\setminus\Delta$. In the tangent
cone this stratum is the flat factor
$\CC^{2n}$. If $p$ is diagonal, an isometry
also preserves the topological nonmanifold
locus of $\Sigma$, so it sends $\Delta$
to the flat factor $\CC^n$ of the diagonal
cone. Lemma~\ref{lem:stratum-length} shows
that in either case the restricted map
would be an isometry of the intrinsic
Riemannian metrics on open subsets of
these strata.

On the local ordered cover of
$\Sigma\setminus\Delta$, the harmonic metric
is $g_N\oplus g_N$. On $\Delta$ it is
$2g_N$, since the tangent inclusion is
$a\mapsto a\otimes I_2$ and
$\tr(I_2)=2$. An isometry to a flat factor
therefore makes $g_N$ flat on a nonempty
open subset of $N$.

The metric $g_N$ is hyperk\"ahler, hence
Ricci-flat and real analytic in harmonic
coordinates \cite[Theorem~5.2]{dk81}.
Analytic continuation of its curvature on
the connected manifold $N$ would make it
flat everywhere. A complete simply
connected flat manifold is Euclidean,
whereas $N$ is compact, simply connected
and of positive dimension. This
contradiction proves the proposition.
\end{proof}

The proposition concerns exact isometry.
It provides no lower bound on the first
nonzero error term in
\eqref{eq:main-comparison}. Optimality of
the exponent, derivative estimates near
lower strata, and extensions to moduli
with a singular-sheaf boundary remain
separate questions.
